\chapter[CAPÍTULO II]{PLAN DEL PROYECTO DE INNOVACIÓN / MEJORA / CREATIVIDAD}

\section{Identificación del Problema Técnico de la Empresa}

\subsection{Problema Local de la Empresa}

En la empresa J\&P Periféricos S.A.C., el Área de Soporte Técnico atiende incidencias, solicitudes de asistencia y requerimientos de mantenimiento de equipos informáticos correspondientes a empresas clientes distribuidas en distintas sedes y sucursales. A partir de la entrevista realizada al Gerente General y al Jefe de Soporte Técnico y TI, junto con la revisión de los registros operativos cotidianos, se evidencian diversas dificultades que afectan la continuidad, el orden y la oportunidad del servicio. Esta situación crítica se desglosa a través de una serie de problemáticas encadenadas:

\begin{enumerate}[label=\alph*)]
    \item Desorganización y falta de trazabilidad en la atención de solicitudes recibidas por canales informales como llamadas telefónicas y mensajes de WhatsApp: Los requerimientos de los clientes corporativos ingresan de manera descentralizada y desorganizada a través de canales informales (llamadas telefónicas y chats de mensajería instantánea), sin un formato estandarizado ni categorización previa, aspecto que, de acuerdo con los lineamientos de gestión operativa expuestos por Gonzales Delgado (2025), propicia que las solicitudes queden dispersas y sin registro oficial en los dispositivos del personal técnico.

    \item Inexistencia de un historial centralizado de mantenimientos y expedientes técnicos por equipo, debido al uso de archivos de Excel dispersos y desactualizados: Aunque el control del parque informático se gestiona parcialmente, existe una total ausencia de convenciones de registro estructurado por empresa, sucursal y trabajador. Se emplean hojas de cálculo descentralizadas y desactualizadas, lo cual, tal como señalan Castro y Quispe (2025), obliga a los operadores a realizar búsquedas manuales que incrementan el riesgo de diagnósticos erróneos y reescritura de informes.

    \item Retrasos y cruces de agenda al coordinar visitas técnicas y atenciones presenciales en las instalaciones de los clientes: Actualmente, coordinar las visitas presenciales y asignar técnicos a campo depende exclusivamente de la comunicación empírica y la memoria operativa, una limitación que, según advierten Ramírez y Gómez (2024), desencadena cruces de agenda y traslados duplicados que perjudican la puntualidad del servicio \textit{in-situ}.

    \item Ausencia de indicadores de gestión (KPIs) y reportes consolidados que permitan medir los tiempos de respuesta y el desempeño del área técnica: No existe un núcleo centralizado que supervise el ciclo de vida de los tickets ni evalúe métricas de rendimiento (KPIs), limitación que, en línea con lo evidenciado por Táez Valdiviezo (2021) al contrastar registros manuales frente a una mesa de ayuda estructurada bajo métricas ITIL 4, impide a la gerencia evaluar objetivamente los tiempos de respuesta y resolución, el cumplimiento de los acuerdos de nivel de servicio (SLA) y la productividad del equipo técnico.

    \item Dificultad para controlar y verificar el cumplimiento de los límites de servicio establecidos en los acuerdos con los clientes (SLA): La verificación de la vigencia de los planes contratados, la cantidad máxima de equipos cubiertos y los tiempos máximos de atención comprometidos se realiza de forma manual y fragmentada; esta carencia de control automatizado, tal como advierten Chávez Gomez (2023) y Figueroa Pozo (2023), dificulta priorizar oportunamente las incidencias críticas, provoca que se brinde soporte a activos fuera de cobertura y compromete el cumplimiento formal de los Acuerdos de Nivel de Servicio (SLA).
\end{enumerate}

\vspace{0.8em}

A partir del diagnóstico aplicado y del análisis de Pareto desarrollado en los Anexos, se determina como problema principal: \textit{Desorganización y falta de trazabilidad centralizada en la gestión de solicitudes de soporte técnico y registro de inventario informático multinivel}. En la práctica, esta situación provoca demoras en la asignación y atención de los casos, dificulta ubicar el historial previo de cada activo antes de intervenirlo, genera reprocesos en el seguimiento de las visitas y afecta la percepción de calidad y cumplimiento por parte de los clientes corporativos.

\vspace{0.8em}

\noindent Con base en el levantamiento de información en el área y los resultados del diagnóstico de Pareto, las causas que originan este problema principal se detallan a continuación:

\begin{enumerate}[label=\alph*)]
    \item \textit{La ausencia de una plataforma web centralizada que reemplace la recepción informal de incidencias por mensajería y llamadas}: Según lo manifestado por la jefatura de soporte técnico durante el diagnóstico interno, la ausencia de una plataforma web unificada que reemplace el uso de mensajería informal impide estructurar adecuadamente el flujo de entrada de los requerimientos, generando registros dispersos que anulan cualquier posibilidad de análisis formal de incidencias.
    
    \item \textit{La falta de una base de datos unificada de equipos, componentes y licencias que funcione como expediente técnico confiable}: Al carecer de una base de datos relacional centralizada y normalizada que funcione como expediente técnico por cliente y sucursal, los operadores recurren a búsquedas manuales propensas a errores, degradando la calidad de la información y dificultando los diagnósticos especializados.
    
    \item \textit{La programación empírica y poco estructurada de las rutas y horarios para las visitas técnicas presenciales}: De acuerdo con la coordinación diaria de campo reportada por el área, se carece de una programación lógica para la automatización de citas y asignación optimizada de personal técnico, limitando la coordinación a métodos manuales que incrementan los tiempos de inactividad en las visitas presenciales.
    
    \item \textit{La carencia de tableros analíticos (KPIs) e indicadores de control que permitan evaluar de forma continua los tiempos de atención y la satisfacción del cliente}: La falta de un sistema con backend robusto y modular que centralice las reglas de negocio impide la monitorización en tiempo real de los indicadores de rendimiento y la validación automatizada de contratos, sobrecargando la operación y comprometiendo la calidad del servicio técnico.
\end{enumerate}

\subsection{Problema Nacional}
A nivel nacional, el ecosistema empresarial peruano está conformado de manera predominante por las Micro, Pequeñas y Medianas Empresas (Mipymes), las cuales constituyen el 99,3\% del tejido empresarial formal y concentran el 89,1\% del empleo privado formal (Ministerio de la Producción [PRODUCE], 2025). A pesar de este rol estructural en la economía, gran parte de estas organizaciones presenta un rezago en su madurez tecnológica y digital; operan bajo esquemas empíricos y continúan dependiendo de canales informales de comunicación (como mensajería instantánea) y registros manuales o hojas de cálculo desactualizadas para gestionar sus actividades cotidianas (PRODUCE, 2025). Esta falta de estandarización limita su capacidad operativa y se ve reflejada directamente en el sector tecnológico y de servicios informáticos —como en la empresa J\&P Periféricos S.A.C.—, donde la ausencia de plataformas de gestión centralizadas impide un adecuado control de incidentes y del inventario de activos TI.

Esta problemática refleja una constante preocupante en el sector tecnológico y de servicios informáticos en el Perú. De acuerdo con Altamira (2026), el proceso acelerado de transformación digital en las micro, pequeñas y medianas empresas (PyMEs) tradicionales exige la migración inmediata de sus flujos operativos hacia plataformas de gestión centralizadas para mantener su competitividad en el mercado nacional. No obstante, un gran porcentaje de organizaciones dedicadas al soporte técnico continúa operando mediante esquemas empíricos y no estructurados, apoyados de forma precaria en aplicaciones de mensajería instantánea y registros manuales (Fiberlux Corp, 2025). Esta falta de estandarización limita la capacidad organizativa de las empresas para escalar sus operaciones, controlar adecuadamente los contratos de servicio y brindar una respuesta oportuna frente a incidentes críticos (Vargas, 2024).

En consonancia con este panorama, reportes técnicos advierten que la carencia de una mesa de ayuda unificada y de un inventario informático estructurado en las organizaciones peruanas impide determinar con precisión la causa raíz de las fallas recurrentes de hardware y software (Castro y Quispe, 2025). Esta limitación operativa incrementa sustancialmente los costos por inactividad empresarial (\textit{downtime}), propicia el cruce constante de agendas técnicas en las atenciones presenciales e imposibilita la verificación de los acuerdos de nivel de servicio (SLA) (Ramírez y Gómez, 2024). Ante esta realidad, el sector tecnológico en el Perú demanda soluciones de software accesibles que integren el control jerárquico de activos, la trazabilidad del ticket en tiempo real y la gestión automatizada de visitas técnicas.

En este contexto de aceleración digital, la adopción de herramientas tecnológicas resulta determinante para garantizar la sostenibilidad y competitividad de las empresas peruanas. Al respecto, Alvarado (2023) sostiene que la transformación digital orientada a la automatización de procesos y la simplificación operativa no solo incrementa la eficiencia y reduce costos, sino que disminuye drásticamente los tiempos de ciclo en la atención de requerimientos del cliente, factor crítico para fortalecer la lealtad e incrementarla satisfacción en el mercado nacional. Sin embargo, la falta de integración de plataformas accesibles y estructuradas en el sector de servicios TI impide que muchas organizaciones locales reduzcan sus tiempos de respuesta, limitando su capacidad para ofrecer un servicio transparente y de calidad global.

Aunado a lo anterior, las deficiencias en la gestión de servicios y el tiempo de inactividad operativo generan pérdidas financieras críticas en el tejido empresarial del país. De acuerdo con un informe especializado publicado por Tricentis (2025), el 40\% de las organizaciones peruanas estima pérdidas anuales que oscilan entre US\$\ 1 millón y US\$\ 5 millones, mientras que otro 42\% reporta pérdidas de entre US\$\ 500 mil y US\$\ 1 millón a causa de fallas de software y la necesidad de mantenimiento correctivo constante. Esto demuestra que la falta de herramientas automatizadas y de una trazabilidad centralizada de incidencias en las empresas locales no solo provoca retrasos operativos, sino que se convierte en un riesgo financiero latente que compromete la continuidad del negocio.


\subsection{Problema Internacional}

A nivel global, la gestión de servicios de tecnología de la información (ITSM) enfrenta desafíos análogos vinculados con la madurez operativa y la automatización de procesos. Según estudios desarrollados por el Service Desk Institute (2024), las mesas de ayuda internacionales priorizan la eliminación paulatina del trabajo manual mediante la incorporación de flujos de trabajo automatizados para la recepción, clasificación, asignación y seguimiento de requerimientos. La dependencia de canales de atención informales es catalogada a nivel internacional como uno de los factores de mayor riesgo para la continuidad del negocio, ya que incrementa la insatisfacción del cliente y genera cuellos de botella en el soporte de primer y segundo nivel (Gartner, 2025).

Por su parte, diversas investigaciones y experiencias internacionales evidencian que operar sin herramientas centralizadas de Help Desk ni métricas formales de servicio dificulta el control de los tiempos de respuesta y resolución de incidencias, mientras que la adopción de plataformas estructuradas bajo métricas ITIL 4 permite reducir dichos tiempos y mejorar el seguimiento operativo (Táez Valdiviezo, 2021). La ausencia de herramientas analíticas centralizadas impide a la gerencia evaluar objetivamente el desempeño del personal técnico y validar el cumplimiento de los contratos de soporte. Asimismo, la adopción de plataformas basadas en buenas prácticas reduce significativamente las horas de inactividad operativa y optimiza la atención de incidencias en organizaciones de diversa índole (Velasco Castro, 2023; Figueroa Pozo, 2023).

En esta misma línea, Díaz Hidalgo (2024) demuestra en su estudio para la Cruz Roja Ecuatoriana que la alta demanda de soporte y el uso de herramientas obsoletas o al límite de su capacidad generan saturación operativa y cuellos de botella en las áreas de tecnología; por ello, la implementación de soluciones Help-Desk en la nube bajo marcos como ITIL permite mitigar estas deficiencias, estabilizando el servicio y optimizando de forma integral los procesos de TI. En consecuencia, la tendencia internacional se orienta hacia la adopción de plataformas web robustas y multitenant que integren en un único entorno la comunicación instantánea, el inventario jerárquico de activos y la analítica de servicios.

A nivel internacional, la gestión ineficiente de las solicitudes y requerimientos en las áreas de tecnología de la información constituye un obstáculo crítico para la productividad institucional. Diversas organizaciones enfrentan desafíos análogos al carecer de herramientas automatizadas para el soporte técnico. Por ejemplo, en la empresa Técnica Avícola S.A. de Pacasmayo (Perú), la ausencia de un sistema de atención informática generaba deficiencias sustanciales en el departamento de TI, lo que motivó la implementación de una plataforma de tipo \textit{Help Desk} para optimizar el registro y la resolución de incidencias (Gonzales, 2016).

Asimismo, en el sector empresarial peruano, la empresa Telvicon evidenció problemas operativos similares relacionados con el manejo desorganizado de sus servicios técnicos, situación que requirió la adopción de un \textit{Service Desk} bajo el marco de trabajo ITIL para mejorar los tiempos de respuesta y la calidad del servicio (Rodríguez y Falcón, 2019). De manera análoga, en el contexto ecuatoriano, la Universidad Politécnica Salesiana analizó el impacto negativo de los métodos tradicionales de soporte, concluyendo que la falta de herramientas centralizadas deteriora el desempeño del área técnica y alarga los tiempos de atención a los usuarios (Ponce y Samaniego, 2015). 

Estas experiencias internacionales demuestran que el uso de procesos manuales o informales en la gestión de incidencias provoca congestión operativa, retrasos y pérdida de información, evidenciando la necesidad imperativa de implementar sistemas web especializados que automaticen el flujo de trabajo en el soporte tecnológico.

\vspace{0.8em}

\subsection{Formulación del Problema}

\noindent \textbf{Problema General}

\vspace{0.5em}

¿De qué manera la implementación de un sistema web Help-Desk centralizado optimiza la gestión de atenciones técnicas y el registro de cambios de hardware y software en equipos multinivel para la empresa J\&P Periféricos S.A.C.?

\vspace{1.2em}

\noindent \textbf{Problemas Específicos}

\begin{itemize}[leftmargin=1.27cm]
    \item ¿De qué forma la asignación automatizada de tickets influye en la reducción de retrasos ocasionados por la gestión manual e informal de solicitudes?
    
    \item ¿Cómo la integración de un módulo de inventario informático por cliente impacta en la trazabilidad del expediente técnico de hardware y software?
    
    \item ¿De qué manera la programación de citas y soporte presencial reduce los cruces de horario y la descoordinación en las visitas \textit{in-situ}?
    
    \item ¿De qué forma el control digital de planes y contratos de soporte contribuye al cumplimiento de los límites de servicio acordados?
    
    \item ¿Cómo el flujo estandarizado de atención favorece el seguimiento y resolución oportuna de los tickets escalados a nivel técnico?
    
    \item ¿De qué manera la provisión de tableros con indicadores de rendimiento (KPIs) permite medir los tiempos de respuesta y la satisfacción del cliente?
\end{itemize}

\newpage

\section{Objetivo del Proyecto de Innovación}

\subsection{Objetivo General}

Desarrollar e implementar un Sistema Web Help-Desk Centralizado para la gestión del servicio de soporte técnico y el control de inventario de activos TI en la empresa J\&P Periféricos S.A.C., con el fin de optimizar los tiempos de atención, garantizar la trazabilidad de los recursos informáticos y mejorar la satisfacción del cliente.

\vspace{0.8em}

\subsection{Objetivos Específicos}

\begin{itemize}[leftmargin=1.27cm]
    \item Automatizar la recepción y asignación de solicitudes de soporte, mediante un módulo de gestión de tickets que elimine los retrasos del canal manual e informal.
    
    \item Implementar un módulo de gestión de expedientes e inventarios de TI (hardware y software), parametrizado por empresa, sucursal, área y trabajador para garantizar la trazabilidad de los equipos.
    
    \item Integrar un módulo de agendamiento y programación de citas presenciales, para optimizar la coordinación y asignación del personal de soporte \textit{in-situ}.
    
    \item Desarrollar un módulo de control digital de planes y contratos de soporte, que valide de forma automatizada los límites de cobertura y Acuerdos de Nivel de Servicio (SLA) vinculados a cada cliente.
    
    \item Estandarizar el procedimiento de derivación y escalamiento de tickets, para asegurar el seguimiento continuo y la resolución oportuna de las atenciones a nivel técnico.
    
    \item Construir un módulo de métricas y tableros de control (KPIs), para evaluar el Tiempo Medio de Resolución (MTTR), medir los tiempos de respuesta y monitorear la satisfacción del cliente.
\end{itemize}

\section{Antecedentes del Proyecto de Innovación / Mejora / Creatividad (investigación realizada)}

\subsection{Antecedentes Nacionales}

En esta tesis de pregrado de la Universidad César Vallejo en la Escuela Profesional de Ingeniería de Sistemas, sustentada por Fonseca Calcin (2025), titulada ``Help desk para la atención de incidencias del área de soporte técnico de una asociación educativa cultural'', se señala como problema principal que la falta de un Help Desk estructurado expone la información institucional a riesgos por virus y amenazas informáticas en la red, dificultando la detección y respuesta oportuna ante incidentes de seguridad, lo que genera desorden, demoras en la solución de fallas y afecta la satisfacción de los usuarios. El objetivo principal fue implementar un sistema de Help Desk para optimizar la gestión de incidencias y reducir los riesgos de la información ante la amenaza de virus e infecciones informáticas en las redes de la institución. La metodología aplicada correspondió a una investigación básica y cuantitativa sobre sistemas de Help Desk y herramientas de gestión tecnológica actuales para resolver problemas de atención de incidencias e inseguridad de la información. Como resultado, se obtuvo que el Help Desk explica el 69.4\% de la mejora en la atención de incidencias, logrando una correlación positiva fuerte que optimiza la detección, registro y respuesta oportuna ante amenazas y fallas en la red. Se concluyó que, a pesar de las constantes amenazas informáticas, la implementación del Help Desk permite combatir los riesgos y proteger la información del servidor mediante una atención de incidencias organizada, rápida y efectiva.

\vspace{0.8em}

En esta tesis de Maestría de la Universidad César Vallejo, sustentada por Villalva Mendivil (2022), titulada ``Sistema Helpdesk en la Gestión de Incidencias del Área de TI en una Empresa de Telecomunicaciones, Lima, 2022'', se expone como problemática que la ausencia de un sistema Helpdesk estructurado generaba desorden en el registro de fallas y demoras prolongadas en la atención de TI, lo que obligaba a realizar desplazamientos innecesarios para problemas menores, afectando el cumplimiento de los tiempos de respuesta (SLA), interrumpiendo la operatividad y reduciendo la satisfacción de los usuarios. El objetivo fue implementar un sistema Helpdesk para mejorar la gestión de incidencias en el área de TI de la empresa TELEGEN PERU SAC mediante un diseño preexperimental. La metodología empleó una investigación aplicada de enfoque cuantitativo y diseño preexperimental de corte longitudinal, evaluando la implementación mediante fichas de registro para medir la mejora en los indicadores de atención, solución en primer nivel, cumplimiento de SLA y tiempos de respuesta. Como resultado, se obtuvo que la implementación incrementó la atención de incidencias del 63.74\% al 88.68\% y la resolución en primer nivel del 48.52\% al 93.57\%, reduciendo sustancialmente los tiempos de respuesta y mejorando de manera significativa el cumplimiento de los SLA. Se concluyó que el sistema Helpdesk optimiza significativamente la gestión de incidencias en el área de TI, permitiendo reducir los tiempos de respuesta y garantizar una atención organizada, rápida y efectiva a través del incremento en la resolución de primer nivel y el cumplimiento de los SLA.

\vspace{0.8em}

En esta tesis de la Universidad Peruana los Andes, presentada por Chávez Cotera (2022), titulada ``Implementación de un Sistema de Help Desk para Mejorar la Gestión de Incidencias del Parque Informático en el Gobierno Regional Junín'', se expone como problemática que en el Gobierno Regional Junín la gestión de incidencias se registraba mediante cuadernos de control manuales o en hojas de cálculo al final de cada semana. Esta falta de automatización generaba una brecha significativa en la atención, atendiéndose solo un promedio del 74\% de los incidentes y asignándose el 80\% de los mismos, lo cual afectaba la calidad de atención y provocaba insatisfacción en los usuarios e ineficiencia en horas-hombre. El objetivo fue determinar la influencia de un sistema Help Desk en la mejora de la gestión de incidencias del parque informático mediante un diseño preexperimental. La metodología empleó una investigación aplicada de nivel explicativo con diseño preexperimental (pre-test y post-test) bajo la metodología RUP, aplicando fichas de registro y un muestreo probabilístico de reportes diarios. Como resultado, se obtuvo que la implementación del sistema Help Desk incrementó el porcentaje de incidencias atendidas del 74.12\% al 74.99\% y el porcentaje de incidencias asignadas del 80.26\% al 83.75\%, mejorando el control y la canalización operativa de los requerimientos. Se concluyó que la implementación del sistema Help Desk optimiza significativamente la gestión de incidencias del parque informático en el Gobierno Regional Junín, superando las limitaciones de los registros manuales y garantizando una atención más organizada y eficiente.

\vspace{0.8em}

En esta tesis de Maestría de la Universidad César Vallejo, sustentada por Figueroa Santander (2022), titulada ``Aplicación ITIL en la Gestión de Incidencias en el Área de Informática de una Empresa de Turismo, Lima 2022'', se expone como problemática que antes de aplicar la metodología ITIL, la empresa presentaba graves deficiencias en la atención de fallas informáticas debido a la falta de un alineamiento adecuado y el seguimiento de procesos. Esto generaba retrasos, cuellos de botella en la recepción de solicitudes, acumulación de incidencias no resueltas y un elevado tiempo de resolución, afectando la calidad del servicio brindado a los usuarios y clientes. El objetivo fue determinar la influencia de la aplicación ITIL en la gestión de incidencias del área de informática de una empresa de turismo en Lima mediante un diseño preexperimental. La metodología correspondió a una investigación aplicada de nivel explicativo y enfoque cuantitativo, con un diseño preexperimental de corte longitudinal. Mediante fichas de observación sobre una muestra probabilística de 215 incidencias (distribuidas en 30 registros), se evaluó el comportamiento de los indicadores de resolución y tiempos antes y después del tratamiento experimental. Como resultado, se obtuvo que la implementación de las buenas prácticas de ITIL incrementó el promedio de incidencias resueltas del 56.52\% al 70.13\% y redujo el porcentaje de incidencias no resueltas del 43.47\% al 29.86\%. Asimismo, el tiempo promedio de resolución de incidencias se redujo sustancialmente de 38.36 a 12.58 minutos. Se concluyó que la aplicación de ITIL influye positivamente en la gestión de incidencias del área de informática, permitiendo optimizar el flujo de trabajo, reducir los tiempos de espera y garantizar una solución organizada, rápida y oportuna a los requerimientos técnicos de la empresa.

\vspace{0.8em}

En esta tesis de Maestría de la Universidad César Vallejo, sustentada por Mendoza Lizana (2023), titulada ``ITIL-V.4 en la Gestión de Incidencias para el área de service desk en una empresa privada de servicios, Lima 2023'', se presenta como problemática que en la entidad de servicios estudiada el proceso de reporte y atención de incidencias resultaba poco fluido, requiriendo que los usuarios enviaran correos a la mesa de ayuda para luego ser evaluados y escalados a distintas áreas para la validación de permisos. La falta de una metodología o flujo estandarizado generaba demoras significativas en la resolución de los requerimientos e incidentes reportados. El objetivo general fue determinar cómo influye ITIL V.4 en la gestión de incidencias para el área de service desk en una empresa privada de servicios en Lima durante el año 2023. La investigación fue de tipo básica, nivel descriptivo-explicativo, con un enfoque cuantitativo y un diseño no experimental, transversal y correlacional, sobre una población censal de 65 colaboradores mediante la técnica de la encuesta. Como resultado, se obtuvo una relación positiva muy alta del 93.7\% entre ITIL V.4 y la gestión de incidencias, influyendo en un 88.9\% en la solución de problemas, un 85.6\% en la calidad del servicio y un 90.1\% en la satisfacción del usuario. Se concluyó que la adopción del marco de buenas prácticas ITIL V.4 influye de manera directa y significativa en la gestión de incidencias del área de service desk, optimizando el flujo de trabajo, reduciendo los tiempos de atención y elevando la calidad del servicio técnico y la satisfacción general de los usuarios.

\vspace{1.2em}

\subsection{Antecedentes Internacionales}

En esta tesis de maestría de la Universidad Técnica del Norte, sustentada por Táez Valdiviezo (2021), titulada ``Implantación de una herramienta cloud Help Desk para mejorar el tiempo de respuesta troubleshooting de red basado en las métricas ITIL4 para el Gobierno Autónomo Descentralizado Municipal del Cantón La Joya de los Sachas'', se señala como problemática que la gestión tradicional de incidencias y soporte técnico —apoyada en registros manuales en hojas de cálculo y comunicación dispersa— generaba retrasos en la atención, tiempos de respuesta y resolución prolongados, así como dificultades para controlar el cumplimiento de los niveles de servicio (SLA). El objetivo general fue implantar una herramienta cloud Help Desk basada en métricas ITIL 4 para mejorar los tiempos de respuesta en el troubleshooting y soporte técnico. La metodología correspondió a un enfoque cuantitativo de tipo aplicado y comparativo (antes y después), contrastando los registros históricos previos (en hojas de cálculo) frente a los tiempos operativos registrados tras la puesta en marcha de la plataforma ServiceDesk Plus en el área administrativa y técnica. Como resultado, se obtuvo una reducción del tiempo promedio de primera respuesta de 10,03 a 4,32 minutos (49,13\% de mejora promedio reportada por el autor) y una disminución del tiempo promedio de resolución de 108,28 a 98,27 minutos (14,46\% de mejora promedio reportada por el autor), además de una mejora en los tiempos por nivel de SLA de 104,07 a 92,40 minutos (11,64\% de mejora promedio reportada por el autor), consolidando un 25,08\% de mejora general promedio en los tiempos de atención. Se concluyó que la adopción de una herramienta Help Desk estructurada bajo métricas ITIL 4 optimiza los tiempos de respuesta y resolución, formaliza el seguimiento de incidencias y brinda mayor control operativo frente a los registros manuales.

\vspace{0.8em}

En esta tesis de la Universidad Politécnica Salesiana, sustentada por Díaz Hidalgo (2024), titulada ``Diseño de software de Help Desk en la nube basado en ITIL para la Cruz Roja Ecuatoriana'', sse identifica como problemática la saturación de la plataforma local de mesa de ayuda (Spiceworks Desktop) debido al alto volumen diario de tickets y al agotamiento de la capacidad y vida útil del servidor donde se alojaba, lo que afectaba el rendimiento y los tiempos de atención del área de TI. El objetivo general fue diseñar e implementar un servicio de Help Desk en la nube basado en las buenas prácticas de ITIL para optimizar recursos y mejorar los tiempos de respuesta del soporte técnico. La metodología aplicó un enfoque inductivo, analítico y experimental bajo modalidad bibliográfica y de campo, sin muestreo poblacional por estar dirigido directamente al departamento de TI de la institución, definiendo roles de prioridad, una fórmula de criticidad de incidentes y la migración del historial de tickets en formato JSON hacia Spiceworks Cloud. Como resultado, se logró migrar el historial conservando los identificadores y comentarios de los tickets, eliminar los costos de servidor físico, licenciamiento y alojamiento en data center, y disponer de un dashboard estadístico para monitorear tickets abiertos/cerrados, categorías y tiempos promedio de primera respuesta y cierre. Se concluyó que el uso de una mesa de ayuda en la nube alineada a ITIL libera infraestructura local, estandariza la atención según la criticidad del incidente y facilita el control operativo de los tiempos de respuesta mediante indicadores visuales.

\vspace{0.8em}

En esta tesis de la Universidad Pública de El Alto, sustentada por Velasco Castro (2023), titulada ``Sistema de Ticket especializado para la Gestión de incidencias en Soporte Técnico'', se expone como problemática que en el gobierno municipal de Achocalla el registro de incidencias se realiza de forma manual mediante formularios preimpresos, ocasionando pérdida de tiempo e información, así como demoras en el registro y atención de los requerimientos. El objetivo general fue desarrollar un sistema de ticket especializado para la gestión de incidencias en soporte técnico que permita el registro automatizado y eficiente de incidencias de hardware y software, mejorando la calidad y rapidez del servicio mediante metodologías ágiles. La investigación adoptó un razonamiento inductivo con enfoque cualitativo, empleando observaciones directas, revisión de formularios preimpresos y entrevistas semiestructuradas, además de validaciones mediante pruebas de caja blanca y negra, métricas ISO/IEC 9126, controles ISO/IEC 27002 y estimación de esfuerzo con el modelo COCOMO II. Como resultado, la implementación demostró un alto nivel de calidad bajo la norma ISO/IEC 9126 con una usabilidad del 95\%, garantizando la confidencialidad, integridad y disponibilidad mediante cifrado y control de accesos por roles. Se concluyó que el sistema web HelpDesk optimiza significativamente la gestión de incidencias en el área de TI, superando las limitaciones de los procesos manuales y garantizando una atención organizada, rápida y efectiva a través del incremento en la eficiencia operativa, el control centralizado de los requerimientos y una usabilidad sobresaliente para los usuarios finales.

\vspace{0.8em}

En esta tesis de pregrado de la Universidad Estatal Península de Santa Elena, sustentada por Figueroa Pozo (2023), titulada ``Estructura de una mesa de ayuda en el departamento de TI del Ministerio de Inclusión Económica y Social con el fin de mejorar el control de incidentes basadas en la metodología ITIL V3'', se expone como problemática que la falta de una mesa de ayuda estructurada y la presencia de un solo encargado para atender a más de 350 usuarios generaban desorden, cuellos de botella y demoras prolongadas en la atención de TI, gestionando las incidencias sin priorización clara ni seguimiento formal al carecer de un catálogo de servicios definido y del marco ITIL v3. El objetivo general was implementar una propuesta de estructuración de la mesa de ayuda basada en las buenas prácticas de ITIL v3 para el departamento de TI del Ministerio de Inclusión Económica y Social, con el fin de optimizar el control de incidentes, reducir los tiempos de atención y mejorar la prestación del servicio. La metodología presentó un enfoque mixto (cuantitativo y cualitativo) de tipo descriptivo y no experimental, empleando encuestas dirigidas al personal administrativo para medir el nivel de satisfacción, entrevistas estructuradas al responsable de TI y observación directa para diagnosticar el flujo de atención y el volumen de incidentes bajo el marco de ITIL v3. Como resultado, se diseñó la estructura de la mesa de ayuda alineada a ITIL v3, definiendo un catálogo de servicios, flujos de atención y niveles de escalado claros; además, la propuesta demostró que, al estandarizar los procesos de recepción y priorización, se eliminan los cuellos de botella, registrando y controlando el 100\% de los incidentes y mejorando los tiempos de respuesta y la satisfacción de los usuarios. Se concluyó que la reestructuración de la mesa de ayuda bajo ITIL v3 optimiza significativamente la gestión de incidentes al estandarizar los procesos de recepción, priorización y escalado, mientras que la incorporación de un catálogo de servicios categorizado y acuerdos de nivel de servicio (SLA) mejora el tiempo de respuesta, elimina el desorden operativo y eleva la calidad percibida del servicio por parte de los usuarios institucionales.

\vspace{0.8em}

En esta tesis de pregrado de la Facultad de Tecnología Franca Dr. Thomaz Novelino, sustentada por Teixeira Ferreira y Knoll Ricardo (2026), titulada ``HelpSynk - Sistema de gerenciamento de tickets: Solução freemium para equipes de TI com atendimento ao cliente'', se plantea como problemática que las pequeñas empresas y equipos internos de TI enfrentan serias dificultades en la gestión de solicitudes debido a la falta de herramientas centralizadas, recurriendo a canales informales (como correo electrónico o mensajería instantánea), lo que genera saturación, duplicación de requerimientos, pérdida de información y una nula trazabilidad en los costos y repuestos utilizados por sector. El objetivo general fue proyectar y desarrollar el sistema web ``HelpSynk'' para la gestión centralizada de tickets de soporte técnico en pequeñas organizaciones, optimizando el flujo de atención, reduciendo los tiempos de respuesta y permitiendo el control de costos operativos por sector. La metodología aplicó un desarrollo de software ágil e incremental, basado en elicitación de requisitos (entrevistas semiestructuradas, encuestas a usuarios finales y análisis comparativo de mercado), además de modelado de procesos en BPMN y diagramas UML (casos de uso, clases, actividades, estados y secuencia). Como resultado, se logró la implementación exitosa de la plataforma web HelpSynk, validada mediante pruebas piloto en un ambiente real con la empresa Hidromar, demostrando una reducción de hasta el 40\% en el tiempo de clasificación de solicitudes, eliminando el registro de llamadas duplicadas y ofreciendo un panel de control interactivo con exportación de reportes a Excel, detallando los costos de repuestos y servicios por atención. Se concluyó que la implementación de un sistema Help Desk propio centraliza de manera eficiente las atenciones informáticas en pequeñas y medianas organizaciones.

\vspace{1.2em}

\subsection{Artículos Científicos}

En el artículo científico publicado por Vargas Calapucha y Romero Licuy (2026), titulado ``Sistema web de gestión de mantenimiento y baja de equipos del Instituto Superior Tecnológico Tena'', se plantea como problemática que el uso de métodos manuales, registros físicos y hojas de cálculo para la gestión de equipos tecnológicos provocaba pérdida de información, demoras en la atención de fallas y serias dificultades en la trazabilidad de los activos. El objetivo principal fue desarrollar una plataforma web para la gestión del mantenimiento y baja de equipos tecnológicos en el Instituto Superior Tecnológico Tena, permitiendo centralizar la información, automatizar procesos y generar reportes para la toma de decisiones. La metodología correspondió a una investigación aplicada con enfoque mixto, empleando el marco ágil SCRUM para el desarrollo iterativo y flexible del software web. Como resultado, se logró la construcción e implementación de la plataforma web, evidenciando una mejora significativa en la organización del mantenimiento, la seguridad de los datos y la eficiencia operativa de los servicios técnicos e institucionales. Se concluyó que el desarrollo del sistema web optimizó los recursos tecnológicos e impulsó la transformación digital del instituto, resolviendo la pérdida de información y garantizando un control eficiente del ciclo de vida de los activos.

\vspace{0.8em}

En el artículo científico de Gonzales Delgado (2025), titulado ``Sistema informático HelpDesk en la calidad de servicio TI de la Dirección Regional de Transportes y Comunicaciones - Tarapoto, 2025'', se analiza la problemática asociada a las limitaciones en la gestión de incidencias registradas antes de la solución (97 no resueltas y 267 horas de tiempo total de resolución), provocando que la calidad del servicio de TI se concentrara en niveles medio (58.2\%) y bajo (41.8\%). El objetivo fue determinar de qué manera el uso de un sistema informático HelpDesk mejora la calidad del servicio de TI en la Dirección Regional de Transportes y Comunicaciones de San Martín -- Tarapoto, 2025. La metodología fue de tipo aplicada con un diseño preexperimental de corte longitudinal; la población y muestra estuvo conformada por 67 usuarios del servicio de TI, empleando la encuesta y la observación mediante cuestionarios y guías de observación. Como resultado, tras implementar el sistema GLPI basado en ITIL 4, las incidencias resueltas aumentaron a 432 (+18\%), las no resueltas cayeron a 32 (-67\%) y el tiempo de resolución bajó a 178 horas (-33\%); asimismo, la calidad alta alcanzó el 52.2\% y se eliminó por completo el nivel bajo (0.0\%). Se concluyó que el uso del sistema informático HelpDesk mejora significativamente la calidad del servicio TI ($p < 0.05$) en todas sus dimensiones (fiabilidad, capacidad de respuesta, seguridad, empatía y elementos tangibles).

\vspace{0.8em}

En el artículo científico publicado por Arteaga y Astudillo León (2012), titulado ``Sistema de Helpdesk mediante el uso de entorno web'', se expone como problemática la falta de un canal web accesible e inmediato para el reporte de incidentes informáticos durante la jornada laboral, sumada a la ausencia de constancia formal en las solicitudes y la falta de un registro estructurado de soluciones técnicas. El objetivo general fue proporcionar un sistema web Helpdesk rápido y sencillo que permita a los empleados solicitar soporte con respaldo vía correo electrónico, consultar la base de conocimiento y calificar la atención recibida. La metodología correspondió a una investigación aplicada de nivel descriptivo, orientada al desarrollo de una plataforma web con integraciones de notificaciones por correo, bitácora de soluciones técnicas y módulos de evaluación de satisfacción. Como resultado, se implementó el sistema web permitiendo el registro centralizado de solicitudes, la emisión de respaldos por correo, la consulta eficiente de la base de conocimiento por parte de soporte y la calificación del nivel de satisfacción por los usuarios. Se concluyó que el sistema de Helpdesk optimiza la atención de incidentes informáticos, garantizando trazabilidad en los requerimientos, reduciendo tiempos mediante la bitácora de soluciones y asegurando la mejora continua del servicio.

\vspace{0.8em}

En el artículo científico publicado por Chávez Gomez (2023), titulado ``Implementación del Sistema de Help Desk Para la Gestión de Incidentes en la Central de Tecnologías de Información de una Entidad del Estado, 2023'', se expone como problemática la desorganización en la atención de solicitudes técnicas, los tiempos elevados de respuesta, la falta de centralización en la comunicación con soporte y las dificultades para dar cumplimiento a los Acuerdos de Nivel de Servicio (SLA). El objetivo general fue implementar un sistema de Help Desk para optimizar la gestión de incidentes, problemas y consultas de usuarios en la Central de Tecnologías de Información de una Entidad del Estado. La metodología correspondió a una investigación aplicada de diseño preexperimental basada en buenas prácticas de ITIL v4, orientada a la estandarización de flujos de trabajo de ticketing y la automatización del soporte técnico. Como resultado, se logró la estandarización de procesos, la recopilación centralizada de datos para análisis detallados, la mejora en la colaboración interna del equipo técnico y una atención más rápida de requerimientos. Se concluyó que el sistema Help Desk contribuye de manera efectiva a la gestión de incidentes, asegurando la calidad en el servicio al cliente, el cumplimiento de los SLA y la mejora continua de los servicios tecnológicos.

\vspace{0.8em}

En el artículo científico publicado por Loza Calderon (2024), titulado ``Revisión de un sistema web helpdesk para la gestión de incidencias de la infraestructura tecnológica en el Hotel GHL Lago Titicaca Puno 2024'', se analiza como problemática las deficiencias en el seguimiento y monitoreo de incidencias del sistema Helpdesk actual del hotel, evidenciadas por la falta de herramientas en tiempo real y la limitada accesibilidad de los usuarios a la información relevante. El objetivo general fue revisar e implementar mejoras en el sistema web Helpdesk del Hotel GHL Lago Titicaca para optimizar la gestión de incidencias tecnológicas, elevar la eficiencia operativa y mejorar la satisfacción del cliente. La metodología adoptó una investigación de enfoque sistemático y aplicado, centrada en el diagnóstico del proceso de soporte técnico, la integración de herramientas analíticas de diagnóstico en tiempo real y la evaluación periódica de los tiempos de respuesta. Como resultado, se incorporaron herramientas de monitoreo en tiempo real, reportes analíticos avanzados y herramientas de diagnóstico, logrando reducir significativamente los tiempos de respuesta/resolución e identificar patrones de fallas para una gestión preventiva. Se concluyó que la optimización del sistema web Helpdesk mejoró la operatividad del hotel y la experiencia de los huéspedes mediante una toma de decisiones más estratégica basada en datos, recomendando evaluaciones periódicas y capacitación continua del personal.

\section{Justificación del Proyecto de Innovación / Mejora / Creatividad}

\subsection{Justificación Legal}
La implementación del sistema HelpDesk para la empresa J\&P Periféricos S.A.C. responde a la necesidad de adecuar las operaciones de soporte técnico al marco normativo peruano sobre gestión y protección de datos. El tratamiento de información de clientes, contactos y credenciales exige el cumplimiento estricto de la \textbf{Ley N.° 29733 (Ley de Protección de Datos Personales)} y su Reglamento (D.S. N.° 003-2013-JUS). La centralización de las incidencias en una base de datos relacional segura elimina el manejo informal de información en aplicaciones de mensajería o correos personales no auditables, reduciendo el riesgo de filtraciones o pérdidas.

Asimismo, la plataforma contempla controles de acceso basados en roles (RBAC), autenticación mediante tokens JWT y registro de auditoría (\textit{logs}), alineándose con las exigencias de la \textbf{Ley N.° 30096 (Ley de Delitos Informáticos)} para prevenir accesos no autorizados, alteración de datos o fugas de información. Por último, la gestión formal de Acuerdos de Nivel de Servicio (SLA) dentro del sistema brinda soporte normativo a los contratos de atención, siguiendo los lineamientos de la norma \textbf{NTP-ISO/IEC 27001} en cuanto a trazabilidad y seguridad en la entrega de servicios de TI.

\subsection{Justificación Social}
En el ámbito organizacional, la adopción de una plataforma estructurada mejora directamente la calidad de vida laboral del equipo de soporte y fortalece la relación con los clientes. En la operatividad actual, la atención empírica genera sobrecarga de trabajo, duplicación de esfuerzos y cruces de horario en las visitas técnicas, lo que incrementa el estrés laboral y dificulta el seguimiento de las solicitudes.

El sistema HelpDesk organiza y prioriza los tickets de atención, facilitando la programación transparente de visitas \textit{in-situ} y soportes remotos. Esto optimiza el uso del tiempo técnico, reduce desplazamientos innecesarios y fomenta un entorno de trabajo más equitativo y ordenado. Por su parte, las empresas cliente obtienen mayor visibilidad sobre el estado de sus requerimientos, lo que construye una relación de confianza mutua basada en la previsibilidad y la calidad del servicio.

\subsection{Justificación Teórica}
El proyecto se fundamenta en la aplicación práctica de marcos de trabajo y estándares internacionales de la Gestión de Servicios de Tecnologías de la Información (ITSM). En el plano operativo, se toman como referencia los flujos de gestión de incidentes, peticiones de servicio y activos de \textbf{ITIL v4}, transformando un proceso que se ejecutaba de forma empírica en una operación estandarizada, medible y trazable.

Desde la perspectiva del desarrollo de software, la propuesta valida la implementación de una arquitectura desasociada moderna. El backend adopta el framework \textit{NestJS} basado en TypeScript con arquitectura modular e inyección de dependencias, respaldado por \textit{MySQL} para datos transaccionales, \textit{MongoDB} para el almacenamiento del historial de chats y \textit{Redis} con \textit{WebSockets} para la comunicación en tiempo real. Esta combinación demuestra teóricamente cómo el uso de arquitecturas escalables y tecnologías web consolidadas mejora la mantenibilidad, estabilidad y rendimiento del software empresarial en procesos de soporte técnico.


\section{Marco Teórico y Conceptual}

\subsection{Fundamento teórico del Proyecto de Innovación / Mejora / Creatividad}

\subsubsection{Sistema Web Help-Desk Centralizado}
\noindent \textbf{Concepto, evolución y rol del help-desk centralizado}\par
Un sistema web Help-Desk centralizado se define como una plataforma informática basada en tecnologías web y distribuidas, diseñada específicamente para automatizar, canalizar y gestionar de forma centralizada la atención de incidentes, solicitudes de servicio y requerimientos organizacionales (Díaz Hidalgo, 2024, pág. 9; Gonzales Delgado, 2026, pág. 12). Esta solución integra interfaces interconectadas que transforman el soporte técnico tradicional en flujos digitales estandarizados (Raya et al., 2014, pág. 18; Lira \& Lira, 2021, pág. 10).

Un sistema help-desk constituye un recurso centralizado de soporte técnico diseñado para asistir a usuarios internos y clientes externos en la resolución de incidencias relacionadas con accesos, fallos de software, hardware o plataformas operativas, actuando como el primer y único punto de contacto (SPOC) para la atención de requerimientos tecnológicos (ServiceNow, 2026). En el ámbito de la gestión de servicios de TI (ITSM), operar como un punto único de contacto facilita la administración de tickets, automatiza tareas cotidianas —como la asignación a agentes o el envío de notificaciones—, permite definir actividades claras en los procesos de negocio y posibilita medir indicadores clave de rendimiento (KPIs) para evaluar el desempeño del área técnica y elevar la productividad y satisfacción del usuario (Al-Hawari y Barham, 2021). Asimismo, como sistemas de mesa de ayuda basados en tickets (ticketing systems), permiten a los usuarios registrar solicitudes en cualquier momento sin depender de horarios restringidos, consolidando distintos canales de comunicación y estructurando el ciclo de vida de cada requerimiento desde su apertura, categorización, priorización y asignación hasta su resolución, escalamiento o cierre (Hampel, 2024).

Históricamente, la atención de soporte en muchas organizaciones operaba bajo esquemas informales o dispersos —como consultas telefónicas directas, llamadas o correos electrónicos personales— que carecían de visibilidad, trazabilidad, opciones de priorización y control sobre la carga laboral, además de dificultar la reutilización del conocimiento generado en incidencias previas (Hampel, 2024; ServiceNow, 2026). La evolución hacia portales web centralizados permite unificar los canales de comunicación en una sola plataforma, estandarizando la recepción mediante un sistema de tickets que registra cada solicitud de principio a fin, facilita su clasificación y escalamiento ordenado, permite consultar y filtrar el historial de interacciones de forma asíncrona y garantiza que ningún incidente quede sin seguimiento (Hampel, 2024; ServiceNow, 2026). A su vez, según las necesidades y el alcance organizacional, las soluciones de help desk y ticketing pueden adoptar modelos internos, externos, híbridos, automatizados o externalizados, así como arquitecturas a medida o plataformas comerciales que varían en usabilidad, capacidad de integración multicanal, automatización, escalabilidad y costos de licenciamiento (Al-Hawari y Barham, 2021; Hampel, 2024; ServiceNow, 2026).

\vspace{0.8em}

\noindent \textbf{Modelo multi-empresa, multinivel y estructura organizacional}\par
Para que un sistema web Help-Desk centralizado opere de manera efectiva al atender a distintas organizaciones y estructuras internas, no basta con concentrar la recepción de solicitudes en un solo portal; es necesario definir un modelo lógico y organizacional que ordene la convivencia de múltiples clientes y la jerarquía interna de cada uno (IBM, s. f.; Niemelä, 2023). Desde la arquitectura de software, el enfoque multiinquilino (\textit{multi-tenant}) establece que una única instancia de la aplicación ---junto con su base de datos o recursos compartidos--- puede prestar servicio a diversos inquilinos u organizaciones clientes que operan dentro de la misma plataforma bajo reglas y privilegios definidos (IBM, s. f.; Niemelä, 2023). Un principio fundamental de este modelo es el aislamiento de los datos y el control de accesos: la información de cada organización debe permanecer separada y protegida frente a otros inquilinos, resguardando la privacidad, confidencialidad y seguridad mediante políticas de autenticación, autorización y control de acceso basado en roles o privilegios mínimos (IBM, s. f.; Niemelä, 2023).

A diferencia de los esquemas de inquilino único (\textit{single-tenant}), donde cada cliente requiere instancias o despliegues aislados de manera independiente, el modelo multiinquilino mejora la mantenibilidad, la eficiencia en el uso de recursos, la escalabilidad y el costo operativo al centralizar actualizaciones y configuraciones, aunque exige equilibrar cuidadosamente el nivel de aislamiento de datos, el rendimiento compartido y la administración de planes o niveles de servicio (IBM, s. f.; Niemelä, 2023). Sobre esta separación entre empresas clientes, el modelo multinivel organiza la estructura interna de cada cuenta para reflejar su realidad operativa y geográfica, articulando niveles jerárquicos como la empresa cliente, sus sucursales o sedes físicas, sus áreas o departamentos internos y los trabajadores o usuarios finales, junto con los planes o contratos de soporte y las zonas geográficas de cobertura (IBM, s. f.). Esta jerarquía permite ubicar con precisión dónde se originan las incidencias, qué activos de hardware y licencias de software están asignados a cada sede o usuario, qué condiciones o Acuerdos de Nivel de Servicio (SLA) aplican y qué permisos corresponden a cada rol operativo (administradores, personal de soporte y clientes solicitantes), apoyándose en la arquitectura relacional del sistema, tokens JWT, control de acceso basado en roles (RBAC) y filtrado según el alcance jerárquico del usuario autenticado.

\begin{itemize}
    \item \textbf{Aislamiento de Datos y Multitenencia:} Garantiza que cada organización cliente opere dentro de la plataforma compartida manteniendo sus registros, usuarios, solicitudes y activos lógicamente separados de otras empresas (IBM, s. f.; Niemelä, 2023). En términos arquitectónicos, este enfoque favorece la mantenibilidad, la escalabilidad y el control de costos operativos frente a despliegues aislados por inquilino, siempre que se resguarden la confidencialidad y la integridad mediante reglas estrictas de aislamiento de datos y control de acceso (Niemelä, 2023). En el sistema, esto se materializa mediante el modelo relacional, la autenticación con tokens JWT y el filtrado automático de consultas según el contexto del inquilino y del usuario autenticado.
    \item \textbf{Jerarquía Organizacional por Empresas, Sucursales y Áreas:} Ordena la estructura interna de cada cuenta mediante niveles jerárquicos que reflejan su operación real ---empresa cliente, sucursales o sedes, áreas o departamentos y trabajadores---, integrando además los planes o contratos de soporte y las zonas geográficas de cobertura (IBM, s. f.). Esta jerarquía permite trazar con exactitud el origen de cada incidencia, asociar los activos de hardware y software correspondientes a cada sede o usuario y aplicar los Acuerdos de Nivel de Servicio (SLA) y condiciones vigentes de forma mantenible y escalable.
    \item \textbf{Roles Operativos y Control de Acceso Jerárquico:} Define los permisos y alcances visibles de cada perfil dentro del sistema ---como administradores, personal de soporte y clientes solicitantes--- bajo principios de control de acceso basado en roles y privilegio mínimo (Niemelä, 2023). Al combinar roles operativos con el alcance jerárquico y organizacional del usuario, se asegura que cada actor acceda únicamente a las funciones, empresas, sedes, áreas y tickets que le corresponden según su responsabilidad dentro de la plataforma.
\end{itemize}

\vspace{0.8em}

\noindent \textbf{Ciclo de vida del ticket, estados y reglas de negocio}\par
El ciclo de vida de un ticket abarca el recorrido integral de una solicitud desde su registro hasta su cierre, mientras que el estado expresa su condición en un momento dado y el flujo de trabajo organiza las etapas y criterios necesarios para su avance (ONES Content Team, 2026). Desde la ingeniería de software, este recorrido puede formalizarse mediante una máquina de estados que define los estados válidos ---como abierto, en progreso, en espera externa, resuelto, cerrado o cancelado---, controla transiciones específicas ---como reanudar desde espera o reabrir un ticket resuelto--- y resguarda los estados terminales frente a cambios no válidos (Faustino, 2026, p. 51). A nivel operativo, diferenciar los flujos según la naturaleza del caso ---como incidentes, requerimientos y problemas--- evita tratar toda solicitud bajo un mismo trámite genérico, permitiendo priorizar la restauración rápida de la disponibilidad del servicio en los incidentes y abordar la causa raíz en los casos recurrentes (Vill\'on Alc\'antara \& Figueroa Vel\'asquez, 2021, pag. 61; Vill\'on Alc\'antara \& Figueroa Vel\'asquez, 2021, pag. 65; Vill\'on Alc\'antara \& Figueroa Vel\'asquez, 2021, pag. 67).

Para garantizar la consistencia, el control y la auditoría de cada solicitud, el ciclo de vida debe regirse por reglas de negocio e invariantes explícitas aplicadas a lo largo del proceso (Faustino, 2026, p. 51; Vill\'on Alc\'antara \& Figueroa Vel\'asquez, 2021, pag. 76; Vill\'on Alc\'antara \& Figueroa Vel\'asquez, 2021, pag. 80). Entre estas reglas y controles destacan:

\begin{itemize}
    \item \textbf{Validación de datos y transiciones:} Exigir información mínima obligatoria al registrar o derivar el caso, verificar transiciones de estado permitidas, impedir alteraciones sobre estados terminales y reajustar la prioridad si durante el seguimiento se evidencia que la evaluación inicial no correspondía (Faustino, 2026, p. 51; Vill\'on Alc\'antara \& Figueroa Vel\'asquez, 2021, pag. 76; Vill\'on Alc\'antara \& Figueroa Vel\'asquez, 2021, pag. 77; Vill\'on Alc\'antara \& Figueroa Vel\'asquez, 2021, pag. 82).
    \item \textbf{Priorización, asignación y propiedad:} Determinar la prioridad a partir del impacto y la urgencia del servicio afectado, mantener un único responsable o grupo solucionador activo en cada momento, documentar el motivo al reasignar o escalar y coordinar la vinculación con problemas cuando ocurren incidentes recurrentes (Faustino, 2026, p. 52; Vill\'on Alc\'antara \& Figueroa Vel\'asquez, 2021, pag. 77; Vill\'on Alc\'antara \& Figueroa Vel\'asquez, 2021, pag. 85; Vill\'on Alc\'antara \& Figueroa Vel\'asquez, 2021, pag. 88).
    \item \textbf{Resolución, conformidad y cierre:} Registrar la solución aplicada o el trabajo temporal (\textit{workaround}), obtener la conformidad del usuario cuando el flujo lo requiera y ejecutar el cierre formal dejando trazabilidad completa del caso (Faustino, 2026, p. 51; Vill\'on Alc\'antara \& Figueroa Vel\'asquez, 2021, pag. 86; Vill\'on Alc\'antara \& Figueroa Vel\'asquez, 2021, pag. 89; Vill\'on Alc\'antara \& Figueroa Vel\'asquez, 2021, pag. 90).
    \item \textbf{Trazabilidad, tiempos y métricas:} Conservar un historial auditable de eventos, cambios de estado, asignaciones y marcas de tiempo para medir volúmenes, tiempos de respuesta y solución, cumplimiento de acuerdos de nivel o tiempos objetivo y recurrencia de casos (Faustino, 2026, p. 51; Vill\'on Alc\'antara \& Figueroa Vel\'asquez, 2021, pag. 69; Vill\'on Alc\'antara \& Figueroa Vel\'asquez, 2021, pag. 70; Vill\'on Alc\'antara \& Figueroa Vel\'asquez, 2021, pag. 100).
\end{itemize}

En un sistema web Help-Desk centralizado, estas reglas articulan además las condiciones propias de operación y servicio ---como planes, cobertura, códigos de seguimiento e historial de actuaciones---, reduciendo ambigüedades, pérdidas de información y cierres inconsistentes durante la atención (Faustino, 2026, p. 51; ONES Content Team, 2026; Vill\'on Alc\'antara \& Figueroa Vel\'asquez, 2021, pag. 67).

\vspace{0.8em}

\noindent \textbf{Asignaci\'on, enrutamiento y control de concurrencia}\par
\indent Dentro de un sistema web Help-Desk centralizado, la asignaci\'on y el enrutamiento constituyen funciones operativas clave para dirigir cada solicitud hacia el especialista o canal adecuado, evitando demoras por distribuci\'on manual y reduciendo los tiempos de resoluci\'on (Wijaya \& Oktavia, 2024, p\'ag. 1831). Aunque la literatura reciente ha explorado enfoques basados en aprendizaje autom\'atico para clasificar y enrutar tickets, los propios estudios evidencian que la asignaci\'on sigue siendo una de las actividades menos desarrolladas frente a la creaci\'on o el monitoreo, por lo que en entornos operativos resulta determinante establecer reglas expl\'icitas basadas en criterios de negocio como la disponibilidad del personal, la carga laboral actual, la modalidad de atenci\'on (remota o presencial) y la cobertura geogr\'afica (Wijaya \& Oktavia, 2024, p\'ags. 1835--1836, 1839--1840). En el contexto del sistema propuesto, esta l\'ogica permite distribuir autom\'aticamente los tickets remotos seg\'un el umbral de espera y la carga activa de los t\'ecnicos, as\'i como enrutar y agrupar citas presenciales \textit{in-situ} por zona geogr\'afica, sucursal y franja horaria para evitar desplazamientos duplicados y cruces de agenda.

\indent Sin embargo, cuando m\'ultiples administradores, t\'ecnicos y usuarios operan simult\'aneamente sobre los mismos tickets o franjas horarias, el acceso concurrente sin mecanismos de control puede provocar anomal\'ias cr\'iticas en los datos, tales como actualizaciones perdidas, lecturas no confirmadas, lecturas irrepetibles o lecturas fantasmas (Arduino \& Alfonzo, 2018, p\'ag. 898). Para preservar la integridad y consistencia de la informaci\'on, las operaciones cr\'iticas deben ejecutarse como transacciones respaldadas por las propiedades ACID (atomicidad, coherencia, aislamiento y durabilidad), combinando estrategias de control de simultaneidad seg\'un la naturaleza de cada operaci\'on: esquemas optimistas basados en validaci\'on antes de confirmar cambios en entornos de baja contenci\'on, y esquemas pesimistas o bloqueos temporales cuando existe riesgo de colisi\'on directa sobre un mismo recurso (Arduino \& Alfonzo, 2018, p\'ags. 897--900). A nivel de aplicaci\'on, este principio se complementa con bloqueos temporales (\textit{soft locks}) y transacciones at\'omicas mediante Prisma ORM sobre MySQL y Redis, asegurando que dos t\'ecnicos no tomen el mismo ticket ni reserven simult\'aneamente la misma franja horaria.

\begin{itemize}
    \item \textbf{Asignaci\'on Automatizada por Disponibilidad y Carga Laboral:} Distribuye las solicitudes entrantes entre los agentes disponibles considerando su carga activa de trabajo y reglas operativas predefinidas, evitando cuellos de botella por asignaci\'on manual y acelerando el inicio de la atenci\'on (Wijaya \& Oktavia, 2024, p\'ags. 1831, 1835--1836). En la plataforma, este mecanismo opera mediante casos de uso automatizados que eval\'uan el umbral de espera de los tickets pendientes y los asignan al t\'ecnico remoto con menor carga operativa.
    \item \textbf{Enrutamiento Operativo y Agrupaci\'on de Visitas In-Situ:} Canaliza los requerimientos hacia soporte remoto o presencial seg\'un la naturaleza de la falla, evaluando la zona geogr\'afica de cobertura, la disponibilidad horaria y las citas previamente agendadas en la misma sucursal. De este modo, el motor de programaci\'on recomienda franjas id\'oneas y agrupa atenciones cercanas para reducir tiempos muertos y traslados innecesarios.
    \item \textbf{Control Transaccional y Bloqueos Temporales (Soft Locks):} Resguarda las operaciones simult\'aneas frente a actualizaciones perdidas e inconsistencias mediante transacciones ACID, aislamiento adecuado y validaci\'on previa a la escritura (Arduino \& Alfonzo, 2018, p\'ags. 897--900). Complementariamente, el uso de bloqueos temporales (\textit{soft locks}) y verificaciones anti-cruces impide que dos operadores modifiquen un mismo ticket o reserven de forma concurrente el mismo horario de visita t\'ecnica.
\end{itemize}

\vspace{0.8em}

\noindent \textbf{Comunicación en tiempo real y mensajería asíncrona}\par
El modelo tradicional de petición--respuesta sobre HTTP y las estrategias de consulta periódica (\textit{polling}) resultan insuficientes para escenarios colaborativos que exigen inmediatez sin sobrecargar el servidor ni forzar la recarga de la interfaz (Gourko, 2024, pp. 4--6). Frente a ello, el protocolo \textit{WebSocket} (RFC 6455) establece un canal bidireccional (\textit{full-duplex}) persistente sobre una única conexión TCP, reduciendo la latencia y permitiendo la transmisión continua de eventos entre cliente y servidor (Gourko, 2024, pp. 4--6, 26--27). Asimismo, frente a alternativas como \textit{Long Polling}---costoso en recursos y complejo de escalar---, \textit{Server-Sent Events} (SSE)---limitado a un flujo unidireccional y a restricciones de conexiones simultáneas por navegador--- o \textit{MQTT}---orientado principalmente a entornos IoT---, los \textit{WebSockets} (y su variante segura \textit{WSS} sobre TLS) ofrecen una comunicación interactiva nativa y eficiente en aplicaciones web modernas (Gourko, 2024, pp. 14--17).

Por otro lado, sostener miles de conexiones persistentes y sincronizar estados en tiempo real introduce retos de concurrencia, consumo de memoria y consistencia de datos entre sesiones o instancias (Gourko, 2024, pp. 10--12, 24--25). Para resolverlo, el patrón \textit{Publicador/Suscriptor} (\textit{Pub/Sub}) desacopla a los emisores de eventos de los receptores mediante canales o tópicos, mientras que un intermediario en memoria como \textit{Redis} actúa como \textit{message broker} de baja latencia para distribuir mensajes de forma asíncrona y coordinar el estado del sistema (Gourko, 2024, pp. 24--25, 29--30). La combinación de \textit{WebSockets} con \textit{Redis Pub/Sub} facilita tanto la entrega instantánea de notificaciones como el escalamiento horizontal de la capa de tiempo real (Gourko, 2024, pp. 26--30).

En el contexto del sistema \textit{Help-Desk} y gestión de activos de TI (\textit{ITAM}), estos mecanismos se aplican de la siguiente manera:

\begin{itemize}
    \item \textbf{Canales bidireccionales persistentes mediante \textit{WebSockets}:} Habilitan el chat integrado de soporte técnico por ticket y la actualización inmediata de estados (como cambios de prioridad, asignación o avance de atención) sin recargar la página, aislando las suscripciones por sala o canal correspondiente a cada ticket (Gourko, 2024, pp. 26--28).
    \item \textbf{Mensajería asíncrona y desacoplamiento con \textit{Redis} (\textit{Pub/Sub}):} Gestionan la difusión asíncrona de eventos operativos---como asignaciones automáticas, alertas de citas \textit{in-situ} o notificaciones en tiempo real---manteniendo la fluidez del sistema y desacoplando los productores de eventos de los consumidores activos (Gourko, 2024, pp. 24--25, 29--30).
    \item \textbf{Persistencia híbrida del historial y trazabilidad de interacciones:} Complementan la velocidad en memoria de \textit{Redis} y la inmediatez del transporte \textit{WebSocket} con el almacenamiento persistente en base de datos documental (\textit{MongoDB}), garantizando la conservación íntegra del historial de conversaciones y la auditoría de cada caso de soporte (Gourko, 2024, pp. 26--30).
\end{itemize}

\vspace{0.8em}

\noindent \textbf{Seguridad web, autenticación y autorización (JWT y RBAC)}\par
La protección de las operaciones y datos sensibles en una aplicación web exige mecanismos robustos de autenticación, gestión segura de sesiones, control de permisos por rol y trazabilidad mediante auditoría (Manda et al., 2026, pp. 1--2; Teixeira Ferreira \& Knoll Ricardo, 2026, p. 2). En el ámbito de las APIs REST, la autenticación basada en \textit{JSON Web Tokens} (JWT) permite mantener un estado de sesión escalable, pero su implementación segura requiere evitar tokens de larga duración sin control y mitigar riesgos como la reutilización de credenciales o el almacenamiento inseguro de \textit{refresh tokens} (Manda et al., 2026, pp. 2--3). Para ello, las arquitecturas modernas combinan \textit{access tokens} de corta vigencia con \textit{refresh tokens} rotativos almacenados mediante funciones de \textit{hashing} adaptativo como \textit{bcrypt}, invalidando sesiones activas ante cambios de contraseña y aplicando validación estricta de entradas y manejo centralizado de errores sin exponer trazas internas en producción (Manda et al., 2026, pp. 1--5).

De manera complementaria, la autorización mediante control de acceso basado en roles (\textit{Role-Based Access Control} -- \textit{RBAC}) asocia permisos a perfiles o roles en lugar de asignarlos de forma aislada a cada usuario, restringiendo el acceso a rutas y operaciones según las responsabilidades de cada actor (Manda et al., 2026, pp. 2--4; Figueroa Pozo, 2023, pp. 23--24). En una arquitectura web modular por capas o limpia---como el \textit{backend} desarrollado en \textit{NestJS} y \textit{TypeScript}---, la lógica de negocio se aísla de los controladores y la persistencia, mientras los \textit{guards} y políticas de autorización validan tanto la identidad del token como el alcance jerárquico del usuario sobre recursos, empresas, sucursales o tickets (Díaz Hidalgo, 2024, p. 9; Manda et al., 2026, pp. 3--4). Asimismo, el registro de auditoría (\textit{audit logs}) documenta eventos críticos y cambios de estado con marcas de tiempo y metadatos estructurados para garantizar la rendición de cuentas y la trazabilidad operativa (Villón Alcántara \& Figueroa Velásquez, 2021, p. 69).

En el contexto del sistema \textit{Help-Desk} y gestión de activos de TI (\textit{ITAM}) para J\&P Periféricos S.A.C., estos lineamientos se concretan en los siguientes pilares:

\begin{itemize}
    \item \textbf{Autenticación con \textit{JWT} y manejo seguro de credenciales:} Verifica la identidad de los usuarios mediante \textit{JSON Web Tokens}, protege las contraseñas con \textit{hashing} mediante \textit{bcrypt} y controla el acceso a las rutas protegidas y \textit{guards} del \textit{backend} en \textit{NestJS}, previniendo accesos no autorizados y manipulación indebida de recursos (Manda et al., 2026, pp. 3--5; Figueroa Pozo, 2023, pp. 23--24).
    \item \textbf{Control de acceso basado en roles (\textit{RBAC}) y alcance multinivel:} Restringe las operaciones y vistas según el rol operativo (administrador, soporte técnico o cliente) y el alcance jerárquico de empresa, sucursal o área, asegurando que cada perfil actúe únicamente sobre la información y acciones permitidas (Manda et al., 2026, pp. 2--4; Teixeira Ferreira \& Knoll Ricardo, 2026, p. 2).
    \item \textbf{Registro de auditoría y trazabilidad de eventos críticos:} Conserva bitácoras estructuradas de eventos relevantes---como cambios de estado en tickets, citas y modificaciones en activos o permisos---respaldando la auditoría, la integridad de los datos y el cumplimiento normativo (Villón Alcántara \& Figueroa Velásquez, 2021, p. 69; Manda et al., 2026, pp. 4--5).
\end{itemize}

\subsubsection{Gestión del Servicio de Soporte Técnico}
\noindent \textbf{Fundamentos de ITSM y buenas prácticas (ITIL v4)}\par
Que se abordara aqui: Cómo se organiza la entrega y co-creación de valor en el soporte técnico mediante flujos estandarizados de incidentes, peticiones de servicio, problemas y escalamiento, superando la atención empírica o informal.

La Gestión de Servicios de Tecnologías de la Información (ITSM) constituye un enfoque estructurado orientado a planificar, entregar, operar y mejorar continuamente los servicios tecnológicos para alinearlos con las necesidades operativas y estratégicas de la organización (Mbeka \& Wausi, 2022). Frente a los esquemas tradicionales de soporte técnico —en los que las áreas de TI operan bajo dinámicas pasivo-reactivas y métodos empíricos que dificultan atender oportunamente la demanda de los usuarios—, los marcos de referencia ITSM establecen un conjunto de lineamientos y buenas prácticas estandarizadas para formalizar el ciclo de vida del servicio y elevar su calidad de entrega.

Dentro de los distintos estándares y marcos de referencia existentes en la industria, Information Technology Infrastructure Library (ITIL) destaca como el marco predominante para la gestión de servicios de TI. La evidencia empírica demuestra que la adopción formal de prácticas estructuradas bajo este enfoque —tales como la gestión de incidentes, la gestión de problemas, la gestión de niveles de servicio y el control de accesos— ejerce una influencia positiva y estadísticamente significativa sobre las dimensiones fundamentales de la calidad del servicio técnico, específicamente en la confiabilidad del desempeño, la capacidad de respuesta oportuna y la seguridad o respaldo brindado al usuario. Asimismo, la efectividad de estas prácticas se ve fortalecida cuando la organización dispone de factores habilitantes clave, como el conocimiento técnico del personal en ITSM y herramientas tecnológicas adecuadas que respalden la operación diaria (Mbeka \& Wausi, 2022).

\noindent \textbf{Gestión de incidentes, soporte remoto y visitas in-situ}\par
Que se abordara aqui: Cómo se estructura la atención de fallas y requerimientos tanto de forma remota como presencial, incluyendo la priorización por impacto y urgencia, el diagnóstico, la derivación a campo y la coordinación de citas sin cruces de agenda.


\noindent \textbf{Acuerdos de Nivel de Servicio (SLA) y planes de soporte}\par
Que se abordara aqui: Cómo se definen, controlan y verifican los compromisos de servicio (vigencia de contratos/planes, cobertura de activos, tiempos máximos comprometidos y límites de atención por cliente).


\noindent \textbf{Indicadores clave de rendimiento (KPIs) y calidad del servicio}\par
Que se abordara aqui: Cómo se mide objetivamente el desempeño operativo y la satisfacción del usuario mediante métricas como MTTA (tiempo medio de respuesta/reconocimiento), MTTR (tiempo medio de reparación/resolución), FCR (resolución al primer contacto) y CSAT (satisfacción del cliente).


\subsubsection{Control de Inventario de Activos TI (ITAM)}
\indent Describe la implementación de un sistema de control centralizado para el parque informático de las empresas clientes, permitiendo el registro auditable y trazable de los componentes de hardware y las licencias de software asignadas en cada locación externa (Chávez Cotera, 2022, pág. 18). Esta gestión centralizada garantiza el control operativo sobre el ciclo de vida de los bienes tecnológicos administrados por J\&P Periféricos S.A.C.

\begin{itemize}
    \item \textbf{Trazabilidad y Historial de Hardware y Software:} Administración continua del estado de conservación, ubicación física y asignación de equipos informáticos en las empresas clientes, permitiendo registrar de forma cronológica el historial detallado de los componentes de hardware y software que estuvieron instalados previamente en cada dispositivo, así como las modificaciones, actualizaciones y reemplazos efectuados a lo largo del tiempo.
    
    \item \textbf{Gestión Estructurada de Datos y Auditoría de Activos:} En el sistema web, la información asociada al inventario corporativo de los clientes y su historial de cambios se gestiona mediante modelos relacionales robustos en MySQL con el soporte de Prisma ORM, garantizando la consistencia transaccional y la trazabilidad exacta de qué recursos tecnológicos poseía un activo en un periodo determinado frente a su configuración actual.
\end{itemize}

\subsection{Conceptos y Términos utilizados}

\begin{description}[font=\bfseries, style=nextline, itemsep=1.2em]

    \item[Ticket]
    Registro digital estructurado que documenta el ciclo de vida completo de un requerimiento de servicio o falla técnica dentro del sistema.

    \item[API REST]
    Interfaz de programación basada en el protocolo HTTP que permite el intercambio de datos en formato JSON entre el cliente frontend y el servidor backend.

    \item[NestJS]
    Framework progresivo de Node.js basado en TypeScript que aplica arquitectura modular, inyección de dependencias y principios SOLID para servicios escalables.

    \item[Vite.js]
    Herramienta de desarrollo y empaquetador web de alta velocidad para la construcción e interacción fluida de interfaces de usuario reactivas.

    \item[MySQL]
    Sistema de gestión de bases de datos relacional (RDBMS) utilizado para la persistencia transaccional y estructurada de los datos centrales del negocio.

    \item[MongoDB]
    Base de datos NoSQL orientada a documentos BSON/JSON utilizada para el almacenamiento flexible del historial del chat de soporte y registros transitorios.

    \item[Redis]
    Almacén de datos en memoria utilizado para la gestión de caché y mensajería en tiempo real mediante el patrón Publicador/Suscriptor (Pub/Sub).

    \item[WebSockets]
    Protocolo de comunicación bidireccional y de baja latencia sobre un único socket TCP para la transmisión de eventos en tiempo real.

    \item[JWT (JSON Web Token)]
    Estándar abierto para la transmisión segura de credenciales y autenticación basada en roles (RBAC).

    \item[ORM (Object-Relational Mapping)]
    Técnica de abstracción que mapea tablas relacionales a objetos de código, implementada mediante Prisma ORM para TypeScript.

    \item[KPI (Key Performance Indicator)]
    Métrica cuantitativa empleada para evaluar el desempeño y eficiencia de la gestión del servicio de soporte (MTTR, MTTA, FCR, CSAT).

    \item[SLA (Service Level Agreement)]
    Acuerdo de Nivel de Servicio que establece los parámetros y tiempos máximos comprometidos para la atención y resolución de requerimientos.

    \item[ITIL v4]
    Marco de referencia internacional de mejores prácticas enfocado en la Gestión de Servicios de Tecnologías de la Información (ITSM).

\end{description}